\section{Revisión de Literatura}

Los estudios relacionados del mercado de combustibles se han centrado en investigar las implicancias de su estructura de mercado, recalcando entre varios factores, la relevancia de la heterogeneidad entre estaciones en las dinámicas y diferenciales de precios \parencite{ECKERT2013,HAUCAP2017,KIHM2016}. Entre los objetos de estudio está la dispersión de precios dentro de un mismo mercado, la cual puede responder al nivel de competencia dentro del mismo \parencite{LEWIS2008} o bien a los costos de búsqueda y transparencia de precios \parencite{SOLIS2025,CHANDRA2011,BYRNE2017,LUCO2019,STIGLER1961,BORENSTEIN1991}. \textcite{LUCO2019} estima de manera causal que la transparencia que ofreció la plataforma Bencinas en Línea en Chile fue mayormente aprovechada por las EDS para vigilar el precio de la competencia, con un uso muy limitado por parte de los consumidores, lo cual acabó incrementando los márgenes. \textcite{SOLIS2025} aporta evidencia que sugiere existencia de costos de busqueda económicamente relevantes, y que, la baja divulgación de la plataforma de Bencinas en Línea podría profundizar los efectos anticompetitivos.

Por otro lado, el nivel competitivo dentro de un mercado se ha aproximado midiendo el número o densidad de competidores en un mercado o zona geográfica \parencite{VANMEERBEECK2003,SEN2003,BERGANTINO2020}. Si bien la cantidad de EDS es un factor fundamental de presión competitiva, \textcite{CASTILLO2018} aporta evidencia para Chile especificando que la mayor presión viene de estaciones independientes\footnote{Estaciones sin bandera, que no pertenecen a las grandes franquicias que controlan la mayoría del mercado}, este trabajo se limita a estimar el efecto de la apertura de nuevas estaciones siguiendo la linea de la literatura más reciente \textcite{FISCHER2025,DAVIS2025,KOH2022,TAYLORMUEHLEGGER2026, BERNARDO2018,CARDOSO2022}. Cuando la apertura no se da por un shock exógeno tales como cambios en la legislación o espacios aptos para apertura, se requiere de un supuesto de exogeneidad en el timing de apertura condicional a variables de control, tal como propone y defiende \textcite{ARCIDIACONO2020} y aplican \textcite{FISCHER2025,TAYLORMUEHLEGGER2026} mediante estimadores de diferencias en diferencias. Dicho supuesto es necesario dado que la decisión de entrada es estratégica, por lo que hay endogeneidad en los mercados y ubicación en que las estaciones deciden abrir \parencite{HOUDE2005}. Adicionalmente, la metodología se ha potenciado gracias a la consolidación nuevas y mejores prácticas en estimaciones de diferencias en diferencias \parencite{CALLAWAYSANTANNA2021,SUNABRAHAM2021,DECHAISEMARTIN2020,BORUSYAK2024,GARDNER2024}.

Por último, dado el componente espacial del mercado, la literatura ha tenido presente los modelos de \textcite{SALOP1979,HOTELLING1929,BRESNAHAN1991} para determinar las hipótesis a testear en lo empírico, siendo \textcite{HOUDE2012} una de las referencias más relevantes en cuanto a las aplicaciones del mercado de combustibles, también destacando \textcite{CLEMENZGUGLER2006,GENAKOS2022} como testeos directos del componente espacial en dicho mercado.